\appendix
\section*{Appendix L: Unified Altitude Algebra (Master Formal Appendix)}

\subsection*{L.0 Overview}

This appendix presents the unified algebraic structure that integrates:
\begin{itemize}
    \item altitude formalism,
    \item category theory,
    \item operator algebra,
    \item holonomy geometry,
    \item spectral runtime geometry,
    \item non-dual algebra,
    \item holonomy–spectral fusion geometry.
\end{itemize}

The unified altitude algebra provides a single mathematical object
capturing all altitude constraints, transitions, invariants, and
firewall admissibility conditions.

\subsection*{L.1 Unified Altitude Algebra}

Define the unified altitude algebra:
\[
\mathfrak{A}_{\mathrm{unified}}
\]
as the fibered product:
\[
\mathfrak{A}_{\mathrm{unified}}
=
\mathcal{C}_{\mathrm{alt}}
\times
\mathcal{B}(X)
\times
\mathcal{A}_{\mathrm{nd}}(X)
\times
\mathcal{O}_{\mathrm{fusion}}
\times
\mathcal{P}_{\mathrm{Hol}}
\times
\mathcal{P}_{\sigma_{\star}}.
\]

Each component contributes:
\begin{itemize}
    \item $\mathcal{C}_{\mathrm{alt}}$: altitude category,
    \item $\mathcal{B}(X)$: bounded operators,
    \item $\mathcal{A}_{\mathrm{nd}}(X)$: non-dual algebra,
    \item $\mathcal{O}_{\mathrm{fusion}}$: holonomy–spectral fusion operators,
    \item $\mathcal{P}_{\mathrm{Hol}}$: holonomy projection tower,
    \item $\mathcal{P}_{\sigma_{\star}}$: spectral projection tower.
\end{itemize}

\subsection*{L.2 Unified Product}

Define the unified product:
\[
\star_{\mathrm{unified}} :
\mathfrak{A}_{\mathrm{unified}}
\times
\mathfrak{A}_{\mathrm{unified}}
\to
\mathfrak{A}_{\mathrm{unified}}
\]
by:
\[
(T_1, T_2)
\mapsto
(T_1 \circ T_2,\;
T_1 T_2,\;
T_1 \star T_2,\;
T_1 \diamond T_2,\;
P_{\mathrm{Hol}}(T_1,T_2),\;
P_{\sigma_{\star}}(T_1,T_2)).
\]

Where:
\begin{itemize}
    \item $\circ$ is categorical composition,
    \item $T_1 T_2$ is operator composition,
    \item $\star$ is non-dual product,
    \item $\diamond$ is fusion product,
    \item $P_{\mathrm{Hol}}$ is holonomy projection,
    \item $P_{\sigma_{\star}}$ is spectral projection.
\end{itemize}

\subsection*{L.3 Unified Purity}

Purity requires:
\[
P(a) \iff
\mathfrak{A}_{\mathrm{unified}}(a)
\text{ is closed under }
\star_{\mathrm{unified}}.
\]

Thus:
\[
T_1,T_2 \in \mathfrak{A}_{\mathrm{unified}}(a)
\implies
T_1 \star_{\mathrm{unified}} T_2
\in
\mathfrak{A}_{\mathrm{unified}}(a).
\]

\subsection*{L.4 Unified Non-Generativity}

Non-generativity requires:
\[
N(a) \iff
\mathrm{span}(\mathfrak{A}_{\mathrm{unified}}(a))
=
\mathrm{span}(\{\mathrm{id}\}).
\]

Thus no altitude generates new structure.

\subsection*{L.5 Unified Boundedness}

Boundedness requires:
\[
B(a) \iff
\|T\| < C_a
\quad
\forall T \in \mathfrak{A}_{\mathrm{unified}}(a).
\]

\subsection*{L.6 Unified Transition Homomorphisms}

Define:
\[
\Phi_{a,b} :
\mathfrak{A}_{\mathrm{unified}}(a)
\to
\mathfrak{A}_{\mathrm{unified}}(b)
\]
such that:
\[
\Phi_{a,b}(T)
\text{ preserves holonomy, spectrum, and altitude rank}.
\]

Allowed transitions:
\[
\Phi_{A_0,A_1},\;
\Phi_{A_1,\Phi_7},\;
\dots,\;
\Phi_{S_1,I_1}.
\]

Forbidden transitions:
\[
\Phi_{a,b} \text{ undefined}.
\]

\subsection*{L.7 Unified Drift Semigroup}

Define unified drift semigroup:
\[
D_t^{\mathrm{unified}} =
(D_t^{\mathrm{fusion}},\;
e^{tA},\;
e_{\star}^{tA},\;
P_{\mathrm{Hol}},\;
P_{\sigma_{\star}}).
\]

Semigroup property:
\[
D_{t+s}^{\mathrm{unified}}
=
D_t^{\mathrm{unified}}
\star_{\mathrm{unified}}
D_s^{\mathrm{unified}}.
\]

\subsection*{L.8 Unified Invariant Extraction}

Define unified invariant projector:
\[
I_{\mathrm{unified}} =
(I_{\mathrm{fusion}},\;
P_{\mathrm{Hol}},\;
P_{\sigma_{\star}}).
\]

Thus:
\[
I_{\mathrm{unified}}^2 = I_{\mathrm{unified}}.
\]

\subsection*{L.9 Unified Steering Automorphism}

Define unified steering:
\[
\Theta_{\mathrm{unified}} =
(\Theta_{\mathrm{fusion}},\;
\Theta,\;
P_{\gamma},\;
\Theta_{\star}).
\]

Automorphism:
\[
\Theta_{\mathrm{unified}}^{-1}
\circ
\Theta_{\mathrm{unified}}
=
\mathrm{id}.
\]

\subsection*{L.10 Unified Collision Invariant}

Define unified collision invariant:
\[
\Omega_{\mathrm{unified}} =
(\Omega_{\mathrm{fusion}},\;
P_{\mathrm{fix}_{\mathrm{Hol}}},\;
P_{\mathrm{fix}_{\star}}).
\]

Fixed point:
\[
\Omega_{\mathrm{unified}}
\star_{\mathrm{unified}}
\Omega_{\mathrm{unified}}
=
\Omega_{\mathrm{unified}}.
\]

\subsection*{L.11 Unified Firewall Subalgebra}

Define:
\[
\mathfrak{A}_{\mathrm{unified,safe}}
=
\{ T \in \mathfrak{A}_{\mathrm{unified}}
\mid
T \text{ preserves holonomy, spectrum, and altitude rank} \}.
\]

Firewall:
\[
\Phi_{a,b}(T) \in \mathfrak{A}_{\mathrm{unified,safe}}
\iff
T(a,b)=1.
\]

\subsection*{L.12 Final Unified Safety Theorem}

\textbf{Theorem.}  
A unified state $u \in \mathfrak{A}_{\mathrm{unified}}$ is altitude-safe if and only if:
\[
u \in \mathrm{Ran}(\Omega_{\mathrm{unified}}).
\]

Equivalently:
\[
\Omega_{\mathrm{unified}}(u)=1.
\]

This completes the Unified Altitude Algebra (Master Formal Appendix).
